\chapter*{Abstract}
\thispagestyle{empty}
SND@LHC è un esperimento di fisica delle particelle, situato al CERN, dedicato allo studio dei neutrini prodotti al Large Hadron Collider.
L'event display attualmente utilizzato dalla collaborazione presenta limiti significativi in termini di interattività e usabilità, fornendo unicamente proiezioni bidimensionali della geometria del rivelatore.

Questa tesi descrive la progettazione e lo sviluppo di un nuovo event display tridimensionale e interattivo, basato su OpenGL, pensato per superare tali limiti.
Dopo un'analisi dei requisiti della collaborazione, vengono illustrate le scelte progettuali e  implementative del sistema, fondate sui principi della computer grafica tridimensionale e su un'architettura modulare ed estendibile.

Il software realizzato consente di esplorare la struttura del rivelatore e gli eventi registrati da prospettive arbitrarie, migliorando sensibilmente la comprensione visiva dei dati rispetto al sistema preesistente, come dimostrato dal confronto su dati sperimentali reali.